
\subsubsection{6.1 Interpretation of Results}

The experiments provided measurements of clusterability in SSL embeddings under the tested conditions (SimCLR and LA-SSL, ResNet-50, 20 epochs, Waterbirds dataset). The observed AMI values (approximately 0.28) were below the 0.4 threshold used to indicate reliable clusterability. Linear separability metrics remained high (AUC $\approx$ 0.98) despite low clusterability scores. This dissociation indicates that linear separability and discrete clusterability are independent properties under these conditions.

These results do not support the hypothesis that standard SSL training with InfoNCE creates discrete, geometrically separable spurious clusters measurable via AMI $\geq$ 0.4 within 20 epochs of training. The hypothesis that LA-SSL operates by dispersing spurious clusters was also not supported; AMI increased rather than decreased.

\subsubsection{6.2 Limitations}

\textbf{Training Duration:} The proof-of-concept experiments used 20 epochs rather than the planned 100 epochs. This was done to validate implementation before committing computational resources to full-scale training. Cluster structure may require longer training to emerge. However, AMI values showed minimal change between epoch 10 and epoch 20, and all three mechanism hypotheses failed under these conditions.

\textbf{Architecture Scope:} Only ResNet-50 was tested. High-capacity models such as ViT-H-14 (tested by Mehta et al., 2022) may exhibit different geometric properties. Clusterability may vary with model capacity.

\textbf{Sample Size for Correlation Analysis:} The M2 test involved n=2 data points (two checkpoints), which is insufficient for reliable statistical correlation estimates. This limits the interpretability of the correlation analysis.

\textbf{Experimental Validation Gap:} The h-e1 hypothesis achieved implementation validation (43/43 tests passing) but experimental validation of intervention efficacy was not completed. The mechanism was implemented but not experimentally tested for efficacy prediction.

\subsubsection{6.3 Alternative Explanations}

The low AMI values despite strong spurious correlation (93\%) in the training data suggest that spurious features may not form discrete clusters under the tested conditions. Two potential explanations are consistent with the observations:

\begin{enumerate}
\item \textbf{Continuous gradients:} Spurious features may form continuous gradients in embedding space rather than discrete density modes. Linear classifiers can identify decision boundaries in continuous feature spaces, which would explain the high linear separability (AUC $\approx$ 0.98) despite low discrete clusterability (AMI $\approx$ 0.28).
\end{enumerate}

\begin{enumerate}
\item \textbf{High-dimensional dispersion:} In 2048-dimensional embedding space, spurious signals may be distributed across many dimensions, reducing the density required for k-means to identify discrete clusters while still permitting linear separability.
\end{enumerate}

Direct visualization or dimensionality reduction analysis was not performed and would be needed to distinguish between these explanations.

\subsubsection{6.4 Implications}

The results indicate that cluster-based diagnostics such as AMI may not identify spurious correlation structure in SSL embeddings under the tested conditions (20-epoch training, ResNet-50). Methods that assume cluster structure exists, such as k-means-based subgroup discovery, may require validation that clusters actually form before being applied.

The finding that linear separability remains high despite low clusterability suggests that linear-based rather than cluster-based diagnostics may be more appropriate for analyzing spurious features in SSL embeddings under these conditions.

The LA-SSL result (AMI increase rather than decrease) indicates that the mechanism by which learning-speed resampling affects fairness may not involve geometric cluster dispersion, at least within 20 epochs of training. Alternative mechanisms involving linear decision boundaries or per-group margins remain possible but were not tested.

\subsubsection{6.5 Future Work}

Full-scale 100-epoch experiments are needed to determine whether cluster structure emerges with extended training. Testing on high-capacity architectures (e.g., ViT-H-14) would establish whether clusterability varies with model capacity. Direct visualization of embedding geometry through dimensionality reduction (t-SNE, UMAP) and analysis of principal components would help characterize the structure of spurious features. Testing cluster-based interventions when high-AMI conditions can be established (if possible) would validate the h-e1 implementation's intervention efficacy predictions. Analysis of per-group linear decision boundaries and margins could test alternative mechanistic explanations for LA-SSL's documented fairness improvements.
